\documentclass[header={Appendix A Exam Study Notes}]{study-note}

\begin{document}

\StudyTitleBlock{A}{Appendix A: Exam Study Notes}{Arzel\`a--Ascoli theorem and compactness in spaces of continuous functions}

\vspace{0.8em}

\begin{focusbox}
This appendix gives one of the main compactness criteria in analysis. In the course it is used especially to prove compactness of integral operators on spaces of continuous functions.
\end{focusbox}

\tableofcontents

\section{Setting}

\begin{defbox}{Space $C(M)$}
Let $M$ be a compact metric space. Then $C(M)$ denotes the space of continuous scalar-valued functions on $M$ with the sup norm
\[
\norm{f}_{\infty}:=\sup_{x\in M}\abs{f(x)}.
\]
\end{defbox}

\begin{defbox}{Equicontinuity}
A family $\mathcal F\subset C(M)$ is \textbf{equicontinuous} if for every $x\in M$ and every $\varepsilon>0$ there exists $\delta>0$ such that
\[
d(x,y)<\delta
\quad\Longrightarrow\quad
\abs{f(x)-f(y)}<\varepsilon
\qquad \forall f\in\mathcal F.
\]
\end{defbox}

\section{Main Theorem}

\begin{thmbox}{Arzel\`a--Ascoli theorem}
Let $M$ be a compact metric space and let $\mathcal F\subset C(M)$. Then $\mathcal F$ is relatively compact in $C(M)$ with the sup norm if and only if:
\begin{enumerate}
\item $\mathcal F$ is bounded in $\norm{\cdot}_{\infty}$;
\item $\mathcal F$ is equicontinuous.
\end{enumerate}
\end{thmbox}

\begin{focusbox}
So in $C(M)$, compactness is controlled by two ingredients: no blow-up in size and no wild oscillation.
\end{focusbox}

\section{Why The Assumptions Are Natural}

\begin{exbox}{Necessity}
If $\mathcal F$ is relatively compact, then its closure is compact. Compact subsets of a normed space are bounded, so $\mathcal F$ is bounded. Equicontinuity follows because a compact family in the uniform norm cannot develop arbitrarily sharp local oscillations.
\end{exbox}

\begin{exbox}{Sufficiency: proof idea}
Assume $\mathcal F$ is bounded and equicontinuous.
\begin{enumerate}
\item Choose a countable dense set $(x_n)$ in $M$.
\item Use diagonal extraction to get pointwise convergence on the dense set.
\item Use equicontinuity and compactness of $M$ to upgrade this to uniform convergence on all of $M$.
\end{enumerate}
Thus every sequence in $\mathcal F$ has a uniformly convergent subsequence, which gives relative compactness.
\end{exbox}

\section{Course Use}

\begin{thmbox}{Compact integral operator on $C([a,b])$}
\thmIntegralOperator
\end{thmbox}

\paragraph{Reason.}
The image of the unit ball is uniformly bounded and equicontinuous, so Arzel\`a--Ascoli applies.

\begin{warningbox}{What to check in applications}
When using Arzel\`a--Ascoli, do not try to prove compactness directly. Instead verify exactly two properties for the image family:
\begin{enumerate}
\item uniform boundedness;
\item equicontinuity.
\end{enumerate}
\end{warningbox}

\section{Exam-Style Exercises}

\begin{exbox}{Exercise A.1: Equicontinuity checklist}
Let
\[
\mathcal F=\{f\in C^1([0,1]): \norm{f}_{\infty}\le 1,\ \norm{f'}_{\infty}\le 1\}.
\]
Prove that $\mathcal F$ is relatively compact in $C([0,1])$.

\paragraph{Solution pattern.}
Uniform boundedness is given. Equicontinuity follows from the mean value estimate
\[
\abs{f(x)-f(y)}\le \abs{x-y}.
\]
Then apply Arzel\`a--Ascoli.
\end{exbox}

\begin{exbox}{Exercise A.2: Bounded but not equicontinuous}
Let
\[
f_n(x)=x^n, \qquad x\in[0,1].
\]
Show that $(f_n)$ is uniformly bounded in $C([0,1])$ but has no uniformly convergent subsequence. Identify which Arzel\`a--Ascoli hypothesis fails.

\paragraph{Solution pattern.}
The sequence is bounded by $1$, but its pointwise limit is discontinuous, so no subsequence can converge uniformly. The missing hypothesis is equicontinuity, which fails near $x=1$.
\end{exbox}

\begin{exbox}{Exercise A.3: Compactness of an integral operator}
Let $k\in C([0,1]^2)$ and define
\[
(Kf)(x)=\int_0^1 k(x,y)f(y)\,dy.
\]
Prove that $K:C([0,1])\to C([0,1])$ is compact using Arzel\`a--Ascoli.

\paragraph{Solution pattern.}
Show that $K(B_1)$ is uniformly bounded and equicontinuous. Boundedness uses $\norm{k}_{\infty}$, while equicontinuity follows from uniform continuity of $k$ on the compact square.
\end{exbox}

\begin{exbox}{Exercise A.4: A compact embedding of H\"older functions}
Let $0<\alpha\le 1$ and consider
\[
\mathcal F=\{f\in C^{0,\alpha}([0,1]): \norm{f}_{\infty}\le 1,\ [f]_{C^{0,\alpha}}\le 1\}.
\]
Prove that $\mathcal F$ is relatively compact in $C([0,1])$.

\paragraph{Solution pattern.}
Uniform boundedness is built in. Equicontinuity follows from
\[
\abs{f(x)-f(y)}\le \abs{x-y}^{\alpha}.
\]
Then apply Arzel\`a--Ascoli.
\end{exbox}

\begin{exbox}{Exercise A.5: Why closedness matters for compactness}
Give an example of a relatively compact subset of $C([0,1])$ that is not compact.

\paragraph{Solution pattern.}
Take the open unit ball inside the one-dimensional subspace of constant functions. Its closure is compact because it lies in a finite-dimensional space, so it is relatively compact, but it is not compact because it is not closed.
\end{exbox}

\section{What To Memorize Precisely}

\begin{focusbox}
Know precisely:
\begin{enumerate}
\item the definition of equicontinuity;
\item the statement of the Arzel\`a--Ascoli theorem;
\item how it is used to prove compactness of integral operators on $C([a,b])$;
\item the slogan: bounded plus equicontinuous gives precompact in $C(M)$.
\end{enumerate}
\end{focusbox}

\end{document}
